\par\noindent\textbf{(1)} 单模光子数 n=0,1,… 的配分函数 z=1/(1−e\textasciicircum{}\{−βhν\})，故平均占有 ̄n=(e\textasciicircum{}\{βhν\}−1)\textasciicircum{}\{−1\}。由对数配分函数或几何分布的二阶矩得 Var(n)=̄n(1+̄n)，体现玻色聚束而非Poisson固定方差。
\par\noindent\textbf{(2)} 频带内 M=g(ν)dν 个独立模，平均能量 ⟨U⟩=Mhν̄n，方差 Var(U)=M(hν)²[̄n+̄n²]=hν⟨U⟩+⟨U⟩²/M。第一项正比能量、体现离散量子噪声，第二项与能量平方成正比，为波动/玻色增强项。
\par\noindent\textbf{(3)} 相对标准差为 √Var(U)/⟨U⟩=√[(1+̄n)/(M̄n)]。在Wien区 ̄n≪1 时相对噪声近似1/√(M̄n)，显现光子计数涨落；Rayleigh–Jeans区 ̄n≫1 时趋1/√M，即即便平均光子数很大，单模热场的相对强度涨落不消失。
\par\noindent\textbf{(4)} 若固定M=1后令T增大，̄n→∞但相对涨落仍趋1；若先令V或频带模式数M→∞，相对波动按 M\textasciicircum{}\{−1/2\} 消失，体现自平均。推导依赖独立自由场模，激光相干态服从Poisson光子数分布，不能直接代入热光的玻色方差。
